\cotefile{156-7}
\scanpage{66}
\mspage{65}

\deleted{On a les \unsure{co} constructions \unsure{de|d'un} $\unsure{L|\mathcal{L}}$ muni de $\parallel$, d'un $\Sigma_L^\circ$, ens des parties de $L$ ``fermées''. Considérons les applications canoniques}

% The opening paragraph and the whole following diagram are cancelled by oblique strokes.
\diagram{\begin{tikzcd}[ampersand replacement=\&, column sep=huge, execute at end picture={\draw ([xshift=-3mm]current bounding box.south west)--([xshift=3mm]current bounding box.north east);}]
\deleted{\Sigma_M^\circ}
  \arrow[r, shift left=1.5ex, "\deleted{A\mapsto A\cap L}"]
\& \deleted{\Sigma_L\overset{\mathrm{def}}{=}\unsure{\mathscr{P}|\mathcal{P}}(L)}
  \arrow[l, shift left=1.5ex, "\deleted{A_S=\{X\in M\mid\operatorname{omb}(X)\subset S\}\mathrel{\leftarrow\!\vert} S}"]
\end{tikzcd}}

Cette \unsure{relation} entre supports dans $\unsure{\mathcal{F}|\mathscr{F}}$, et supports dans $M$, \unsure{marcherait} en remplaçant $M$ par un \unsure{sous-ensemble} $\Lambda\subset\mathcal{F}$, soumis à la seule condition que voici :

\noindent\begin{tabular}{|p{0.93\linewidth}}
$\underline{\mathrm{Disj}}_\Lambda$ : Soit $F\in\mathcal{F}$, $\operatorname{omb}_\Lambda(F)=\{X\in\Lambda\mid X\ll F\}$.
Alors pour tout $G\in\mathcal{F}$, on a
\[
F\parallel G\quad\Longleftrightarrow\quad\forall X\in\operatorname{omb}_\Lambda(F),\quad X\parallel G
\]
\end{tabular}

Posant, pour $A\subset\mathcal{F}$
\[
\operatorname{omb}_\Lambda(A)\overset{\mathrm{def}}{=}\bigcup_{F\in A}\operatorname{omb}_\Lambda(F)\subset\Lambda,
\]
on voit bien que pour tte partie $A\in\mathcal{F}$, on a
\[
\operatorname{cosupp}_{\mathcal{F}}A=\operatorname{cosupp}_{\mathcal{F}}\operatorname{omb}_\Lambda A.
\]

De plus, si $A$ est fermé dans $\deleted{\unsure{F}}\mathcal{F}$ pour $\ll$, alors on voit aussitôt que
\[
\operatorname{omb}_\Lambda(A)=\Lambda\cap A\qquad(A\text{ fermé ds }\mathcal{F},\ll)
\]
(en général on a seulement $\supset$). Donc on conclut
\[
\operatorname{cosupp}_{\mathcal{F}}(A)=\operatorname{cosupp}_{\mathcal{F}}A\cap\Lambda
\qquad\begin{aligned}&\text{si }A\text{ fermé}\\&\text{dans }\mathcal{F}\text{ pour }\ll\end{aligned}
\]

Ceci implique, en \unsure{reprenant} les supports
\[
\deleted{\unsure{\Sigma_A|A}}\operatorname{supp}_{\mathcal{F}}(A)=\deleted{\unsure{\mathrm{co}}}\operatorname{supp}_{\mathcal{F}}(A\cap\Lambda)
\qquad\begin{aligned}&\text{si }A\text{ fermé}\\&\text{dans }\mathcal{F}\text{ pour }\ll\end{aligned}
\]
et sans hyp. sur $A$
\[
\operatorname{supp}_{\mathcal{F}}(A)=\operatorname{supp}_{\mathcal{F}}(\operatorname{omb}_\Lambda(A))
\]
